\section{A parallel with proton irradiation}\label{sec:memon}

Memon et al.~\cite{memon2025tr,memon2025arxiv} exposed three commercial
compute modules running Linux to 20 to 58\,MeV protons at the AIC-144
cyclotron of IFJ PAN, Krak\'ow: a Raspberry Pi Zero~2~W, an NXP
i.MX~8M~Plus board (Yocto Linux 5.15.77, root filesystem ext4 on eMMC)
and an OrangeCrab FPGA board running a VexRiscv soft core. They record
133 Linux single-event functional interrupts across the three, and on
the i.MX~8M~Plus attribute 56\,\% of failures to the filesystem and
34\,\% to the storage driver, up to ext4 journal aborts and read-only
remounts~\cite{memon2025arxiv}. The work follows failures of unmodified
Linux on commercial modules down to the kernel subsystem, across three
architectures, with failure cross-sections per platform and energy. That
resolution is what makes it usable by a filesystem designer, and its
finding that the storage path dominates on the i.MX~8M~Plus is the
observation \bfs answers. Their result and ours meet at the storage
path; \cref{tab:memon} sets the two methods side by side.

\begin{table}[t]
\centering
\caption{Physical irradiation~\cite{memon2025arxiv} and this report.}
\label{tab:memon}
\footnotesize
\begin{tabularx}{\columnwidth}{@{}>{\raggedright\arraybackslash}p{1.3cm}
  >{\raggedright\arraybackslash}X >{\raggedright\arraybackslash}X@{}}
\toprule
 & Memon et al. & this report \\
\midrule
cause & protons, 20 to 58\,MeV & emulated single-bit flips \\
where & SoC die (and memory, see text) & read-bio buffers of chosen blocks \\
effects & any: crashes, hangs, driver and filesystem errors &
  wrong bytes delivered to the filesystem only \\
filesystem & ext4 on eMMC & \bfs, ext4, ext3, btrfs, squashfs \\
observable & failure class per run & outcome per read; kernel
  corrections per decode \\
dose & fluence per run & flips received per attack \\
hardware & exposed; replaced by an unexposed unit after each fluence;
  not reported retested & not exposed; medium verified unaltered after
  every attack \\
\bottomrule
\end{tabularx}
\end{table}

The two are complementary. A beam reproduces physics no emulator does:
controller hangs, crashes, the faults their failure classes record. It
is costly, and each unit absorbs dose. Emulation injects one mechanism,
at chosen blocks, with a record of every flip, and leaves the hardware
untouched; it cannot say what a beam does to a controller. \bfs
addresses only the bytes a read returns: against a hung eMMC controller
it can at best fail closed.

\subsection{What a failure class says about the hardware}

A beam acts on hardware; software only reports what follows. Telling
the two apart is what lets a failure be read as a single-event effect on
a sound device rather than as the symptom of a damaged one, so we read
both texts on that point.

Their protocol is careful about cumulative damage. Failures are observed
through software-visible channels, the kernel log, the system journal
and user-space signals, logged off the device, and root causes are
attributed by kernel log forensics; each run ends at its first
functional interrupt~\cite{memon2025arxiv}. Supply current and voltage
are logged throughout. After each targeted fluence every device is
replaced by an identical, previously unexposed copy, so that total
ionising dose and displacement damage do not bias the
analysis~\cite{memon2025arxiv}; the journal manuscript replaces devices
after each test objective and mentions real-time detection of current
latch-ups~\cite{memon2025tr}, while the arXiv text places single-event
latch-up, micro-latch-up included, outside its scope. For one memory
error the authors note that a prior clean run weakens the hypothesis of
a permanent cell defect~\cite{memon2025arxiv}.

What a failure class cannot say is what the beam did to the hardware. On
the i.MX~8M~Plus a filesystem error has at least three physical
readings, which the logs alone do not separate:
\begin{enumerate}[label=(\roman*)]
\item a transient upset in the system on chip, in a core, a cache, the
  storage host controller or a DMA engine, that corrupts data or kernel
  state in flight and leaves the hardware intact;
\item data altered on the eMMC itself, by an upset in the eMMC if it was
  in the beam, or by the system on chip writing corrupted data, which
  survives a power cycle;
\item a part degraded by its dose, through ionising dose, displacement
  damage or a micro-latch-up, that no longer behaves as a fresh one.
\end{enumerate}
They call for different remedies, and only the first is a property of
software running on sound hardware; the second and third are effects on
hardware that software then reports. The checks that separate them are a readback of the storage
against its pre-exposure image, a functional and electrical test of each
unit after its exposure (supply current at rest against its pre-exposure
value, a memory test), and the dose each unit received. Neither text
reports them, and neither states where the units came from. Replacing
units between fluences keeps damage from carrying into the next run; it
does not show that a unit was intact at the end of its own run, and a
cross-section computed over all the runs at one energy assumes every
unit equivalent and unaltered throughout. The modules are commercial parts, not hardened
ones, and their own tolerance is precisely what is left unmeasured: the
campaign establishes how often these systems fail under protons, not how
much the hardware itself withstands, nor, for a given failure, whether
software met a transient upset or a damaged device.

The beam geometry on the i.MX~8M~Plus is also stated three ways: the
journal manuscript puts the system on chip and the LPDDR4 memory in the
beam; the arXiv text puts only the system-on-chip die in the beam
aperture, external memory and storage outside; its Table~1 lists the
system on chip, the LPDDR4 and the eMMC as irradiated. Whether the eMMC
and the memory were exposed changes which readings are possible for
90\,\% of that board's failures.

This matters to \bfs directly. It covers reading (i) when the upset
corrupts the bytes a read returns, and reading (ii) when the altered
bytes stay within its correction capacity; it can do nothing for
corrupted kernel state or a degraded controller, where it can at best
fail closed. Which reading dominates is therefore the first question a
beam campaign on \bfs must answer, and the one this emulation sidesteps
by construction: here the hardware is not exposed, and the medium was
read back intact after every attack (\cref{sec:multifs}).

These are reporting points, not evidence against the measurements, which
we take as they are reported. A beam campaign on \bfs would need them
settled from the start: unit provenance and pre-exposure
characterisation, dose per unit, a post-exposure functional and
electrical test of every unit, a readback of the storage against its
pre-exposure image, and the exact parts in the beam.

\subsection{The same questions, asked of this report}\label{sec:selfaudit}

The checks we ask of a beam campaign apply to this report too, with the
instruments in the place of the hardware. Each item says what was
checked and what was found.

\begin{enumerate}[label=(\alph*)]
\item \emph{What was measured.} The units here are software, named by
  their hashes: the \bfs commit, the Yocto layer, the root image and the
  kernel of each architecture, and \srcTrees{} compiled source trees found
  identical to the sources the signed manifests list (\apprepro). \bfs
  is built into both kernels; the module file its recipe also installs
  is never loaded. The tools are identified by SHA-256 against the
  repository, not by the version they print (mkfs.beamfs prints
  \xfsXMkfsVersion, fsck.beamfs \xfsXFsckVersion, the module is 0.1.26).
  The two kernels differ in one \bfs option: the x86-64 kernel has
  the tree checker \texttt{CONFIG\_BEAMFS\_DEBUG\_TREE}
  (\cref{sec:setup}); the aarch64 kernel, on which every attack ran, does
  not. The x86-64 xfstests sweep ran with more internal checks than the
  attacks.
\item \emph{State after exposure.} The attacks leave the medium
  untouched by construction, and it was read back after every attack:
  twelve files intact on every filesystem (\cref{sec:multifs}). For the
  x86-64 sweep, the node ran without a restart: the harness recorded no
  recovery, and the node's uptime grew by the length of the sweep. Its
  kernel taint flags read \xfsXtaintStart{} at the start and the same at
  the end, without the warning or the oops flag: no kernel warning and
  no oops occurred. The harness clears the kernel log before each test
  and keeps it after; none of the \xfsXkernLogs{} logs holds a warning or
  an oops, and in \xfsXtreeLogs{} of them the tree checker reported no
  lost pointer.
\item \emph{Dose.} Every flip is in the injector's ring and is matched
  against the kernel's own record of each correction
  (\cref{sec:injector}); doses still differ between filesystems
  (\cref{sec:limits}).
\item \emph{What the instruments reported.} Our own tools misreported,
  and each defect was found by reading a raw record against an
  independent account of the same event.
  \begin{itemize}
  \item emufi logs a hook one bio early on this kernel, found against
    the kernel's correction lines (\cref{sec:injector}).
  \item The xfstests harness, before 2.4.0, saved a test that xfstests
    declined as passed, because check counts it in its ``Passed all''
    line; found against check's own output for each test. Before 2.3.61
    the per-test output of a sweep was overwritten by the next one. 2.4.0
    asked the node for a module source version that \bfs does not
    expose, and recorded an empty field. All three are corrected in
    2.5.0, and the x86-64 sweep was run again under 2.5.1; the counts of
    \cref{tab:xfs} are read from check's output, not from the harness's
    summary.
  \item The first x86-64 sweep under 2.5.0 ran under the harness's
    default budget of 1\,900\,s per test, below the \xfsXsoloSeconds\,s
    generic/476 had taken alone on the same image, and the budget, not
    \bfs, stopped the test after \xfsSupFullSeconds\,s. That budget was
    ours to set, and the record did not name it. The sweep was stopped
    after \xfsSupDone{} tests; the harness now records its budget and
    every recovery of the node (2.5.1), and the sweep was run again
    under 14\,400\,s, in which generic/476 passed.
  \item The bench's record of each btrfs attack has no checksum line in
    its slice of the kernel log; the cause of the refusal, a
    \texttt{csum failed} warning on mirror~1, is read from the node's
    own kernel log, where it appears in each run. The bench's topology
    report calls \bfs a loaded module while it is built in. Correction
    lines are rate-limited (\cref{sec:limits}).
  \end{itemize}
\item \emph{What the tools did not reach.} squashfs was never attacked:
  its attack printed no target in any run. The write path received
  no flip. xfstests declined \xfsXNotrun{} tests on x86-64 and
  \xfsANotrun{} on aarch64 (\cref{sec:xfstests}); they are untested.
\end{enumerate}

Two of these defects would have changed a result had the raw records
not been read: which block each flip hit, and the xfstests counts; a
third, the budget, would have reported generic/476 as failed on x86-64.
The
others change no number in this report; they are listed so that they
are corrected before the next campaign. All are in the tools; none is
in the \bfs code measured here.
